\documentclass[../../数学分析培优讲义(答案版).tex]{subfiles}

\begin{document}

\setcounter{chapter}{4}
\pagestyle{fancy}

\chapter{不定积分}

\section{不定积分计算方法}

\subsection{基本方法}

\begin{example}
求下列不定积分:
\begin{enumerate}[label=(\arabic*)]
    \item $\displaystyle I=\int\dfrac{\mathrm{d}x}{3x^2+3x-1}$;
    \item $\displaystyle I=\int\dfrac{\mathrm{d}x}{\sqrt{3x^2+3x-1}}$;
    \item $\displaystyle I=\int\dfrac{\mathrm{d}x}{x^2-5x+7}$.
\end{enumerate}
\end{example}

\begin{proof}[解]
配方后直接使用基本积分公式，得到
\begin{enumerate}[label=(\arabic*)]
    \item $\displaystyle I=\dfrac{1}{\sqrt{21}}\ln\left|\dfrac{6x+3-\sqrt{21}}{6x+3+\sqrt{21}}\right|+C$；
    \item $\displaystyle I=\dfrac{1}{\sqrt3}\ln\left|2\sqrt3\sqrt{3x^2+3x-1}+6x+3\right|+C$；
    \item $\displaystyle I=\dfrac{2}{\sqrt3}\arctan\dfrac{2x-5}{\sqrt3}+C$。
\end{enumerate}
\end{proof}

\begin{example}
求下列不定积分:
\begin{enumerate}[label=(\arabic*)]
    \item $\displaystyle I=\int\dfrac{a^x-a^{-x}}{a^x+a^{-x}}\,\mathrm{d}x\quad(a>0)$;
    \item $\displaystyle I=\int\dfrac{\ln x+1}{1+x\ln x}\,\mathrm{d}x$;
    \item $\displaystyle I=\int\dfrac{\sin2x}{a^2\sin^2x+b^2\cos^2x}\,\mathrm{d}x\quad(a^2\neq b^2)$;
    \item $\displaystyle I=\int(3x-1)\sqrt{3x^2-2x+7}\,\mathrm{d}x$.
\end{enumerate}
\end{example}

\begin{proof}[解]
均将分子配成分母或根式内函数的导数：
\begin{enumerate}[label=(\arabic*)]
    \item 当 $a\neq1$ 时，
    \[
        I=\dfrac{1}{\ln a}\ln(a^x+a^{-x})+C;
    \]
    当 $a=1$ 时，$I=C$。
    \item $\displaystyle I=\ln|1+x\ln x|+C$。
    \item $\displaystyle I=\dfrac{1}{a^2-b^2}\ln|a^2\sin^2x+b^2\cos^2x|+C$。
    \item 令 $u=3x^2-2x+7$，则
    \[
        I=\dfrac{1}{3}(3x^2-2x+7)^{\frac32}+C.
    \]
\end{enumerate}
\end{proof}

\begin{example}
求下列不定积分:
\begin{enumerate}[label=(\arabic*)]
    \item $\displaystyle I=\int(\tan x-\sec x)^2\,\mathrm{d}x$;
    \item $\displaystyle I=\int\sin^4x\,\mathrm{d}x$;
    \item $\displaystyle I=\int\dfrac{\sin^4x}{\cos^2x}\,\mathrm{d}x$;
    \item $\displaystyle I=\int\sin mx\cos nx\,\mathrm{d}x$.
\end{enumerate}
\end{example}

\begin{proof}[解]
利用恒等变形与积化和差公式：
\begin{enumerate}[label=(\arabic*)]
    \item
    \[
        (\tan x-\sec x)^2=2\sec^2x-1-2\sec x\tan x,
    \]
    故 $I=2\tan x-x-2\sec x+C$。
    \item $\displaystyle I=\dfrac{3x}{8}-\dfrac{\sin2x}{4}+\dfrac{\sin4x}{32}+C$。
    \item 由 $\sin^4x/\cos^2x=\sec^2x-2+\cos^2x$，
    \[
        I=\tan x-\dfrac{3x}{2}+\dfrac{\sin2x}{4}+C.
    \]
    \item 当 $m\neq n$ 且 $m\neq-n$ 时，
    \[
        I=-\dfrac{\cos((m-n)x)}{2(m-n)}
        -\dfrac{\cos((m+n)x)}{2(m+n)}+C.
    \]
    当 $m=\pm n\neq0$ 时，按原被积函数直接得到
    $I=-\cos(2mx)/(4m)+C$。
\end{enumerate}
\end{proof}

\begin{example}
求下列不定积分:
\begin{enumerate}[label=(\arabic*)]
    \item $\displaystyle I=\int\dfrac{\mathrm{d}x}{x^4+x}$;
    \item $\displaystyle I=\int\dfrac{x\,\mathrm{d}x}{\sqrt{1-4x^2}}$;
    \item $\displaystyle I=\int\dfrac{\mathrm{d}x}{x^4\sqrt{1+x^2}}$;
    \item $\displaystyle I=\int\dfrac{\mathrm{d}x}{x\sqrt{x^2-1}}$.
\end{enumerate}
\end{example}

\begin{proof}[解]
\begin{enumerate}[label=(\arabic*)]
    \item 由
    \[
        \dfrac{1}{x(x^3+1)}=\dfrac{1}{x}-\dfrac{x^2}{x^3+1},
    \]
    得 $\displaystyle I=\ln|x|-\dfrac13\ln|x^3+1|+C$。
    \item 令 $u=1-4x^2$，得 $\displaystyle I=-\dfrac14\sqrt{1-4x^2}+C$。
    \item 令 $y=\sqrt{1+x^2}/x$，则
    \[
        \mathrm{d}y=-\dfrac{\mathrm{d}x}{x^2\sqrt{1+x^2}},
        \qquad
        \dfrac{1}{x^2}=y^2-1.
    \]
    因而
    \[
        I=y-\dfrac{y^3}{3}+C
        =\dfrac{\sqrt{1+x^2}}{x}
        -\dfrac13\left(\dfrac{\sqrt{1+x^2}}{x}\right)^3+C.
    \]
    \item $\displaystyle I=\arccos\dfrac{1}{|x|}+C\quad(|x|>1)$。
\end{enumerate}
\end{proof}

\begin{example}
求下列函数的不定积分:
\begin{enumerate}[label=(\arabic*)]
    \item $\displaystyle I=\int\dfrac{\mathrm{d}x}{x\ln x}$;
    \item $\displaystyle I=\int\dfrac{\mathrm{d}x}{x^8(1+x^2)}$.
\end{enumerate}
\end{example}

\begin{proof}[解]
\begin{enumerate}[label=(\arabic*)]
    \item $\displaystyle I=\ln|\ln x|+C$。
    \item 利用
    \[
        \dfrac{1}{x^8(1+x^2)}
        =\dfrac{1}{x^8}-\dfrac{1}{x^6}+\dfrac{1}{x^4}-\dfrac{1}{x^2}+\dfrac{1}{1+x^2},
    \]
    得
    \[
        I=-\dfrac{1}{7x^7}+\dfrac{1}{5x^5}-\dfrac{1}{3x^3}
        +\dfrac{1}{x}+\arctan x+C.
    \]
\end{enumerate}
\end{proof}

\begin{example}
求下列不定积分:
\begin{enumerate}[label=(\arabic*)]
    \item $\displaystyle I=\int\dfrac{\mathrm{d}x}{\sin x}$;
    \item $\displaystyle I=\int\dfrac{\mathrm{d}x}{\cos x}$;
    \item $\displaystyle I=\int\dfrac{\cos x}{\sqrt{\cos2x}}\,\mathrm{d}x$;
    \item $\displaystyle I=\int\dfrac{\mathrm{d}x}{3\cos^2x+4\sin^2x}$;
    \item $\displaystyle I=\int\dfrac{\sin x}{1+\sin x}\,\mathrm{d}x$;
    \item $\displaystyle I=\int\dfrac{\mathrm{d}x}{2+\cos^2x}$.
\end{enumerate}
\end{example}

\begin{proof}[解]
\begin{enumerate}[label=(\arabic*)]
    \item $\displaystyle I=\ln\left|\tan\dfrac{x}{2}\right|+C$。
    \item $\displaystyle I=\ln|\sec x+\tan x|+C$。
    \item 令 $u=\sin x$，得
    \[
        I=\dfrac{1}{\sqrt2}\arcsin(\sqrt2\sin x)+C.
    \]
    \item 令 $t=\tan x$，得
    \[
        I=\int\dfrac{\mathrm{d}t}{3+4t^2}
        =\dfrac{1}{2\sqrt3}\arctan\dfrac{2\tan x}{\sqrt3}+C.
    \]
    \item 由 $1/(1+\sin x)=\sec^2x-\sec x\tan x$，
    \[
        I=x-\tan x+\sec x+C.
    \]
    \item 令 $t=\tan x$，得
    \[
        I=\dfrac{1}{\sqrt6}\arctan\left(\sqrt{\dfrac23}\tan x\right)+C.
    \]
\end{enumerate}
\end{proof}

\begin{example}
求下列不定积分:
\begin{enumerate}[label=(\arabic*)]
    \item $\displaystyle I=\int\ln x\,\mathrm{d}x$;
    \item $\displaystyle I=\int xe^x\,\mathrm{d}x$;
    \item $\displaystyle I=\int\ln(1+\sqrt{x})\,\mathrm{d}x$;
    \item $\displaystyle I=\int\left(\dfrac{\ln x}{x}\right)^2\,\mathrm{d}x$.
\end{enumerate}
\end{example}

\begin{proof}[解]
分部积分可得
\begin{enumerate}[label=(\arabic*)]
    \item $\displaystyle I=x\ln x-x+C$；
    \item $\displaystyle I=(x-1)e^x+C$；
    \item 令 $t=\sqrt{x}$，则
    \[
        I=(x-1)\ln(1+\sqrt{x})+\sqrt{x}-\dfrac{x}{2}+C;
    \]
    \item $\displaystyle I=-\dfrac{(\ln x)^2+2\ln x+2}{x}+C$。
\end{enumerate}
\end{proof}

\begin{example}
求下列函数的不定积分:
\begin{enumerate}[label=(\arabic*)]
    \item $\displaystyle I=\int x\cos x\,\mathrm{d}x$;
    \item $\displaystyle I=\int(2x+3x^2)\arctan x\,\mathrm{d}x$;
    \item $\displaystyle I=\int\dfrac{x\arcsin x}{\sqrt{1-x^2}}\,\mathrm{d}x$;
    \item $\displaystyle I=\int\dfrac{\ln\sin x}{\sin^2x}\,\mathrm{d}x$;
    \item $\displaystyle I=\int\sin(\ln x)\,\mathrm{d}x$.
\end{enumerate}
\end{example}

\begin{proof}[解]
逐项分部积分：
\begin{enumerate}[label=(\arabic*)]
    \item $\displaystyle I=x\sin x+\cos x+C$。
    \item 取 $\mathrm{d}v=(2x+3x^2)\,\mathrm{d}x$，得到
    \[
        I=(x^2+x^3+1)\arctan x-x-\dfrac{x^2}{2}
        +\dfrac12\ln(1+x^2)+C.
    \]
    \item $\displaystyle I=x-\sqrt{1-x^2}\arcsin x+C$。
    \item $\displaystyle I=-\cot x\ln|\sin x|-\cot x-x+C$。
    \item 令 $t=\ln x$，则
    \[
        I=\dfrac{x}{2}[\sin(\ln x)-\cos(\ln x)]+C.
    \]
\end{enumerate}
\end{proof}

\subsection{特殊方法}

\methodsection{(A) 配对法}

\begin{example}
求下列不定积分:
\begin{enumerate}[label=(\arabic*)]
    \item $\displaystyle\int\dfrac{\sin x}{\sin x+\cos x}\,\mathrm{d}x$;
    \item $\displaystyle\int\dfrac{\mathrm{d}x}{3+5\tan x}$;
    \item $\displaystyle\int\dfrac{\mathrm{d}x}{\sin x+2\cos x+3}$.
\end{enumerate}
\end{example}

\begin{proof}[解]
\begin{enumerate}[label=(\arabic*)]
    \item 将本积分与分子为 $\cos x$ 的积分配对，得到
    \[
        I=\dfrac{x}{2}-\dfrac12\ln|\sin x+\cos x|+C.
    \]
    \item 设 $D=3\cos x+5\sin x$。由
    \[
        \cos x=\dfrac{5}{34}D'+\dfrac{3}{34}D,
    \]
    可得
    \[
        I=\dfrac{5}{34}\ln|3\cos x+5\sin x|+\dfrac{3x}{34}+C.
    \]
    \item 令 $t=\tan(x/2)$，则
    \[
        I=2\int\dfrac{\mathrm{d}t}{(t+1)^2+4}
        =\arctan\dfrac{t+1}{2}+C.
    \]
\end{enumerate}
\end{proof}

\begin{example}
利用配对法求下列不定积分:
\begin{enumerate}[label=(\arabic*)]
    \item $\displaystyle\int\dfrac{\mathrm{d}x}{1+x^4}$;
    \item $\displaystyle\int\dfrac{\mathrm{d}x}{1+x^2+x^4}$;
    \item $\displaystyle\int\dfrac{\mathrm{d}x}{1+x^3}$.
\end{enumerate}
\end{example}

\begin{proof}[解]
分解实二次因式后可得
\begin{enumerate}[label=(\arabic*)]
    \item
    \[
        I=\dfrac{1}{4\sqrt2}\ln\dfrac{x^2+\sqrt2x+1}{x^2-\sqrt2x+1}
        +\dfrac{1}{2\sqrt2}\arctan\dfrac{\sqrt2x}{1-x^2}+C;
    \]
    \item
    \[
        I=\dfrac14\ln\dfrac{x^2+x+1}{x^2-x+1}
        +\dfrac{1}{\sqrt3}\left[
        \arctan\dfrac{2x+1}{\sqrt3}
        +\arctan\dfrac{2x-1}{\sqrt3}\right]+C;
    \]
    \item
    \[
        I=\dfrac13\ln|x+1|-\dfrac16\ln(x^2-x+1)
        +\dfrac{1}{\sqrt3}\arctan\dfrac{2x-1}{\sqrt3}+C.
    \]
\end{enumerate}
反正切函数跨越分支点时，积分常数相应调整。
\end{proof}

\methodsection{(B) 递推法}

\begin{example}
求下列不定积分的递推公式:
\begin{enumerate}[label=(\arabic*)]
    \item $\displaystyle I_n=\int\tan^n x\,\mathrm{d}x$;
    \item $\displaystyle I_n=\int\sec^n x\,\mathrm{d}x$;
    \item $\displaystyle I_n=\int\dfrac{\sin nx}{\sin x}\,\mathrm{d}x$.
\end{enumerate}
\end{example}

\begin{proof}[解]
\begin{enumerate}[label=(\arabic*)]
    \item 由 $\tan^n x=\tan^{n-2}x(\sec^2x-1)$，
    \[
        I_n=\dfrac{\tan^{n-1}x}{n-1}-I_{n-2},
        \qquad n>1.
    \]
    \item 分部积分得
    \[
        I_n=\dfrac{\sec^{n-2}x\tan x}{n-1}
        +\dfrac{n-2}{n-1}I_{n-2},
        \qquad n>1.
    \]
    \item 由
    \[
        \dfrac{\sin nx-\sin((n-2)x)}{\sin x}=2\cos((n-1)x),
    \]
    得
    \[
        I_n=I_{n-2}+\dfrac{2\sin((n-1)x)}{n-1},
        \qquad n>1.
    \]
\end{enumerate}
\end{proof}

\begin{example}
求下列不定积分的递推公式:
\begin{enumerate}[label=(\arabic*)]
    \item $\displaystyle I_n=\int\dfrac{\mathrm{d}x}{(1+x^2)^n}$;
    \item $\displaystyle I_n=\int\ln^n x\,\mathrm{d}x$;
    \item $\displaystyle I_{m,n}=\int x^m\ln^n x\,\mathrm{d}x$.
\end{enumerate}
\end{example}

\begin{proof}[解]
\begin{enumerate}[label=(\arabic*)]
    \item 分部积分可得
    \[
        I_n=\dfrac{x}{2(n-1)(1+x^2)^{n-1}}
        +\dfrac{2n-3}{2(n-1)}I_{n-1},
        \qquad n>1.
    \]
    \item
    \[
        I_n=x\ln^n x-nI_{n-1}.
    \]
    \item 当 $m\neq-1$ 时，
    \[
        I_{m,n}=\dfrac{x^{m+1}\ln^n x}{m+1}
        -\dfrac{n}{m+1}I_{m,n-1}.
    \]
    当 $m=-1$ 时，
    \[
        I_{-1,n}=\dfrac{\ln^{n+1}x}{n+1}+C.
    \]
\end{enumerate}
\end{proof}

\newpage

\section{几类可积函数}

\subsection{有理分式}

\begin{example}
求下列不定积分:
\begin{enumerate}[label=(\arabic*)]
    \item $\displaystyle\int\dfrac{\mathrm{d}x}{x(x^3+1)^2}$;
    \item $\displaystyle\int\dfrac{x(x^2+3)}{(x^2-1)(x^2+1)^2}\,\mathrm{d}x$;
    \item $\displaystyle\int\dfrac{2x^4+5x^2-2}{2x^3-x-1}\,\mathrm{d}x$;
    \item $\displaystyle\int\dfrac{\mathrm{d}x}{\cos^7x}$.
\end{enumerate}
\end{example}

\begin{proof}[解]
\begin{enumerate}[label=(\arabic*)]
    \item 令 $u=x^3$，则
    \[
        I=\dfrac13\int\dfrac{\mathrm{d}u}{u(u+1)^2}
        =\dfrac13\left[\ln\left|\dfrac{u}{u+1}\right|+\dfrac{1}{u+1}\right]+C.
    \]
    \item 令 $u=x^2$，部分分式分解给出
    \[
        I=\dfrac12\left[\ln\left|\dfrac{u-1}{u+1}\right|+\dfrac{1}{u+1}\right]+C.
    \]
    \item 长除法与部分分式分解给出
    \[
        \dfrac{2x^4+5x^2-2}{2x^3-x-1}
        =x+\dfrac{1}{x-1}+\dfrac{4x+3}{2x^2+2x+1},
    \]
    因此
    \[
        I=\dfrac{x^2}{2}+\ln|x-1|+\ln(2x^2+2x+1)
        +\arctan(2x+1)+C.
    \]
    \item 连续使用 $\sec^n x$ 的递推式，得到
    \[
        I=\dfrac16\sec^5x\tan x+\dfrac{5}{24}\sec^3x\tan x
        +\dfrac{5}{16}\sec x\tan x
        +\dfrac{5}{16}\ln|\sec x+\tan x|+C.
    \]
\end{enumerate}
\end{proof}

\subsection{无理函数}

\begin{example}
求下列不定积分:
\begin{enumerate}[label=(\arabic*)]
    \item $\displaystyle\int\dfrac{x+\sqrt[3]{x^2}+6\sqrt{x}}{x(1+\sqrt[3]{x})}\,\mathrm{d}x$;
    \item $\displaystyle\int\dfrac{\mathrm{d}x}{\sqrt{x+1}+\sqrt[3]{x+1}}$;
    \item $\displaystyle\int\sqrt[3]{\dfrac{2-x}{2+x}}\dfrac{\mathrm{d}x}{(2-x)^2}$;
    \item $\displaystyle\int\dfrac{1}{x}\sqrt{\dfrac{x+2}{x-2}}\,\mathrm{d}x$.
\end{enumerate}
\end{example}

\begin{proof}[解]
\begin{enumerate}[label=(\arabic*)]
    \item 令 $t=\sqrt[6]{x}$，化简后
    \[
        I=6\int\left(t^3+6-\dfrac{6}{1+t^2}\right)\mathrm{d}t
        =\dfrac32t^4+36t-36\arctan t+C.
    \]
    \item 令 $t=\sqrt[6]{x+1}$，则
    \[
        I=6\int\dfrac{t^3}{t+1}\,\mathrm{d}t
        =2t^3-3t^2+6t-6\ln|t+1|+C.
    \]
    \item 令
    \[
        t=\sqrt[3]{\dfrac{2-x}{2+x}},
    \]
    则被积式化为 $-3\mathrm{d}t/(4t^3)$，故
    \[
        I=\dfrac{3}{8t^2}+C
        =\dfrac38\left(\dfrac{2+x}{2-x}\right)^{\frac23}+C.
    \]
    \item 令 $t=\sqrt{(x+2)/(x-2)}$，则
    \[
        I=-2\int\left(\dfrac{1}{t^2-1}+\dfrac{1}{t^2+1}\right)\mathrm{d}t
        =\ln\left|\dfrac{t+1}{t-1}\right|-2\arctan t+C.
    \]
\end{enumerate}
\end{proof}

\begin{example}
求下列不定积分:
\begin{enumerate}[label=(\arabic*)]
    \item $\displaystyle\int\dfrac{\sqrt{x^2+2x+3}}{x}\,\mathrm{d}x$;
    \item $\displaystyle\int\dfrac{\mathrm{d}x}{x^2\sqrt{4x^2-3x-1}}$;
    \item $\displaystyle\int\dfrac{1-\sqrt{1+x+x^2}}{x\sqrt{1+x+x^2}}\,\mathrm{d}x$;
    \item $\displaystyle\int\dfrac{\mathrm{d}x}{\sqrt{x^2+3x-4}}$.
\end{enumerate}
\end{example}

\begin{proof}[解]
\begin{enumerate}[label=(\arabic*)]
    \item 使用 Euler 代换
    \[
        t=\dfrac{\sqrt{x^2+2x+3}+\sqrt3}{x},
        \qquad
        x=\dfrac{2(\sqrt3t+1)}{t^2-1}.
    \]
    化为有理积分后得到
    \[
        I=-\sqrt3\ln|3t+\sqrt3|+\ln\left|\dfrac{t+1}{t-1}\right|
        +\dfrac{1-\sqrt3}{t+1}+\dfrac{1+\sqrt3}{t-1}+C.
    \]
    \item 令 $u=1/x$。在固定符号区间内，积分化为
    \[
        -\int\dfrac{u\,\mathrm{d}u}{\sqrt{4-3u-u^2}},
    \]
    从而
    \[
        I=\sqrt{4-3u-u^2}+\dfrac32\arcsin\dfrac{2u+3}{5}+C.
    \]
    \item 令 $u=1/x$，并合并对数项，得到
    \[
        I=-\ln\left|2\sqrt{1+x+x^2}+x+2\right|+C.
    \]
    \item
    \[
        I=\ln\left|2\sqrt{x^2+3x-4}+2x+3\right|+C.
    \]
\end{enumerate}
\end{proof}

\begin{example}
求下列不定积分:
\begin{enumerate}[label=(\arabic*)]
    \item $\displaystyle\int x^5(a^3+x^3)^{\frac{1}{3}}\,\mathrm{d}x$;
    \item $\displaystyle\int\dfrac{(x^6+a^6)^3}{x^{11}}\,\mathrm{d}x$.
\end{enumerate}
\end{example}

\begin{proof}[解]
\begin{enumerate}[label=(\arabic*)]
    \item 令 $u=x^3$，再令 $v=a^3+u$，得
    \[
        I=\dfrac17(a^3+x^3)^{\frac73}
        -\dfrac{a^3}{4}(a^3+x^3)^{\frac43}+C.
    \]
    \item 展开后逐项积分：
    \[
        I=\dfrac{x^8}{8}+\dfrac{3a^6x^2}{2}
        -\dfrac{3a^{12}}{4x^4}-\dfrac{a^{18}}{10x^{10}}+C.
    \]
\end{enumerate}
\end{proof}

\subsection{三角函数有理式}

\begin{example}
求下列不定积分:
\begin{enumerate}[label=(\arabic*)]
    \item $\displaystyle\int\dfrac{\mathrm{d}x}{a+b\tan x}$;
    \item $\displaystyle\int\dfrac{\mathrm{d}x}{\sin x\cos2x}$;
    \item $\displaystyle\int\dfrac{\cos2x}{\sin^4x+\cos^4x}\,\mathrm{d}x$;
    \item $\displaystyle\int\dfrac{\mathrm{d}x}{1-\sin x}$;
    \item $\displaystyle\int\dfrac{\sin x}{1+\sin x}\,\mathrm{d}x$;
    \item $\displaystyle\int\dfrac{\mathrm{d}x}{5+3\sin x+4\cos x}$;
    \item $\displaystyle\int\dfrac{\sin x}{1+\sin x+\cos x}\,\mathrm{d}x$.
\end{enumerate}
\end{example}

\begin{proof}[解]
\begin{enumerate}[label=(\arabic*)]
    \item
    \[
        I=\dfrac{ax+b\ln|a\cos x+b\sin x|}{a^2+b^2}+C,
        \qquad a^2+b^2\neq0.
    \]
    \item 由
    \[
        \dfrac{1}{\sin x\cos2x}=\dfrac{1}{\sin x}+\dfrac{2\sin x}{\cos2x},
    \]
    得
    \[
        I=\ln\left|\tan\dfrac{x}{2}\right|
        +\dfrac{1}{\sqrt2}\ln\left|\dfrac{\sqrt2\cos x+1}{\sqrt2\cos x-1}\right|+C.
    \]
    \item 令 $t=\tan x$，则
    \[
        I=\int\dfrac{1-t^2}{1+t^4}\,\mathrm{d}t
        =\dfrac{1}{2\sqrt2}\ln\dfrac{t^2+\sqrt2t+1}{t^2-\sqrt2t+1}+C.
    \]
    \item $\displaystyle I=\tan x+\sec x+C$。
    \item $\displaystyle I=x-\tan x+\sec x+C$。
    \item 令 $t=\tan(x/2)$，则
    \[
        I=2\int\dfrac{\mathrm{d}t}{(t+3)^2}
        =-\dfrac{2}{t+3}+C.
    \]
    \item 同样令 $t=\tan(x/2)$，或直接验证，得到
    \[
        I=\dfrac{x}{2}-\dfrac12\ln|1+\sin x|+C.
    \]
\end{enumerate}
\end{proof}

\end{document}

